% ============================================================
% ch10_error_hierarchy.tex
% Error Measurement Hierarchy
% Part of: Tetrad-Based Departure Decomposition (Jiwon Park)
% Session S9 draft — 2026-03-30
% ============================================================

\chapter{Error Measurement Hierarchy}
\label{ch:error-hierarchy}

Our goal in this chapter is to define a systematic error
budget for the species-resolved Bianchi Boltzmann solver of
Chapter~\ref{ch:results} and the multi-species $\Teff$
reduction of Chapter~\ref{ch:teff}.  The budget comprises five
levels, ordered from the microscopic state space to the
macroscopic observable, with each level measuring a distinct
source of error that can be diagnosed and controlled
independently.

The five-level structure addresses a problem that is unique to
the $\Teff$ programme.  A standard Boltzmann solver (CAMB,
CLASS) computes the CMB power spectrum from a single code path
and reports a single numerical convergence criterion
($\ell_{\rm max}$, integration tolerance).  The $\Teff$
solver, by contrast, introduces additional approximation
layers---the effective-temperature ansatz, the species-resolved
tangency diagnostic, the $\ell$-band routing, and the
nonlinear source bridge---each with its own error
characteristic.  Without a layered error budget, it is
impossible to determine where the dominant error arises and
where computational resources should be concentrated.

We define the five levels in
Sections~\ref{sec:err-state}--\ref{sec:err-observable},
discuss the cost--accuracy trade-off in
Section~\ref{sec:err-cost}, and summarise the error budget
for the Phase~1.0 solver in
Section~\ref{sec:err-phase10-budget}.


% ============================================================
\section{Level~1: State-space error}
\label{sec:err-state}
% ============================================================

The first level measures the deviation of the solver's
internal state---the PSTF multipoles of each species---from
a reference solution.


\paragraph{Definition.}

For species $s$ at conformal time $\eta$, angular multipole
$\ell$, azimuthal index $m$, and Fourier wavenumber $k$, the
state-space error is
\begin{equation}\label{eq:err-state}
\varepsilon_{\rm state}^{(s)}(\ell, \eta)
:= \frac{
  \sum_{m,k}
  \big|
    U_{\ell m k}^{(s),\text{Teff}}
    - U_{\ell m k}^{(s),\text{ref}}
  \big|^2
}{
  \sum_{m,k}
  \big|
    U_{\ell m k}^{(s),\text{ref}}
  \big|^2
}\,,
\end{equation}
where $U_{\ell m k}^{(s)}$ denotes the species multipole
(temperature, E-mode, B-mode for photons; energy and number
moments for neutrinos) and the superscripts ``Teff'' and
``ref'' denote the $\Teff$-backed solver and the reference
solution, respectively.


\paragraph{Reference solutions.}

The choice of reference depends on the background:
in the FLRW or near-isotropic limit, the reference is
the standard CAMB or CLASS output; on an anisotropic
Bianchi background, the reference is the same background
evolved with a non-$\Teff$ PSTF hierarchy (i.e.\ the full
multipole hierarchy truncated at $\ell_{\rm max}$ without
the effective-temperature ansatz).  This two-tier baseline
is essential: comparing the $\Teff$ solver against CAMB on
a strongly anisotropic background would conflate the
anisotropy-induced deviation with the $\Teff$
approximation error.


\paragraph{Band-averaged variants.}

For diagnostic purposes, we define band-averaged
state-space errors:
\begin{equation}\label{eq:err-state-band}
\varepsilon_{\rm low}^{(s)}
:= \max_{\ell \leq 2}\,
  \varepsilon_{\rm state}^{(s)}(\ell, \eta_*)\,,
\qquad
\varepsilon_{\rm mid}^{(s)}
:= \max_{3 \leq \ell \leq \ell_{\rm trans}}\,
  \varepsilon_{\rm state}^{(s)}(\ell, \eta_*)\,,
\qquad
\varepsilon_{\rm high}^{(s)}
:= \max_{\ell > \ell_{\rm trans}}\,
  \varepsilon_{\rm state}^{(s)}(\ell, \eta_*)\,,
\end{equation}
where $\ell_{\rm trans}$ is the Silk damping scale
(Section~\ref{sec:ell-band}).  The low-$\ell$ error is
the most relevant for the departure parameter pipeline,
which depends only on $D_2$ ($\ell = 2$).


\paragraph{Phase~1.0 status.}

The Phase~1.0 solver operates at $\ell_{\rm max} = 2$, so
only $\varepsilon_{\rm low}$ is defined.  In the FLRW
limit, $\varepsilon_{\rm low} = 0$ by construction (the
solver recovers $D_2 = 0$ exactly).  On the Bianchi~I
background, the oracle comparison gives
$\varepsilon_{\rm state}^{(\nu)}(2, \eta_*) \simeq 0.08$
(the $-8\%$ deviation of
Section~\ref{sec:disc-oracle}), dominated by the
analytic transfer function's missing late ISW and
neutrino stress feedback.  Phase~2.0 targets
$\varepsilon_{\rm low} < 0.02$ through self-consistent
Poisson (Section~\ref{sec:phase20}).


% ============================================================
\section{Level~2: Source-translation error}
\label{sec:err-source}
% ============================================================

The second level measures the accuracy of the $\Teff$
nonlinear source bridge---the map from the angular
temperature field $\Theta_s(\hat{e})$ to the Thomson
scattering source that enters the Boltzmann hierarchy.


\paragraph{Definition.}

The source-translation error for the photon Thomson source is
\begin{equation}\label{eq:err-Th}
\varepsilon_{\rm Th}
:= \frac{
  \big|
    I_{ab}^{\rm ref}
    - c_\xi\,
      \langle \Theta_\gamma^4\,
      \hat{e}_{\langle a}\hat{e}_{b\rangle}
      \rangle_\Omega
  \big|
}{
  \big| I_{ab}^{\rm ref} \big|
}\,,
\end{equation}
where $I_{ab}^{\rm ref}$ is the photon anisotropic stress
computed from the full Boltzmann hierarchy (without the
$\Teff$ ansatz) and the $\Theta^4$ term is the $\Teff$
moment-map reconstruction
(Eq.~\ref{eq:exact-bridge}).  A companion diagnostic
compares against the linearised bridge:
\begin{equation}\label{eq:err-Th-lin}
\varepsilon_{\rm Th}^{\rm lin}
:= \frac{
  \big|
    I_{ab}^{\rm ref}
    - 4\,c_\xi\,\Theta^{(\gamma)}_{ab}
  \big|
}{
  \big| I_{ab}^{\rm ref} \big|
}\,.
\end{equation}
The difference
$\varepsilon_{\rm Th}^{\rm lin} - \varepsilon_{\rm Th}$
isolates the contribution of the nonlinear ($\Theta^4$
vs.\ $4\Theta$) correction.  At moderate dipole amplitudes
($A \sim 10^{-2}$), the linear bridge can err by up to
$52\%$, while the nonlinear bridge maintains
$\varepsilon_{\rm Th} < 1\%$---a demonstration that the
$\Theta^4$ moment map provides a qualitatively superior
source translation even when the dipole is not small.


\paragraph{Phase~1.0 status.}

In the Phase~1.0 solver, the Thomson source uses the
linearised form (Eq.~\ref{eq:Thomson-damping}) because the
shear-driven quadrupole is the dominant source and the
dipole is small ($A \sim \beta \sim 10^{-3}$).  The
source-translation error is therefore
$\varepsilon_{\rm Th}^{\rm lin} \sim 10^{-3}$, consistent
with the second-order estimate of
Section~\ref{sec:quadratic-sources}.  The nonlinear bridge
becomes essential only when the dipole amplitude exceeds
$\sim\!10^{-2}$, which lies in the nonperturbative regime
(Section~\ref{sec:nonperturbative}).


% ============================================================
\section{Level~3: Adequacy certificate}
\label{sec:err-adequacy}
% ============================================================

The third level is not an error per se but a
\emph{certificate}: a bound on the observable-level
representation error, computed from the tangency diagnostic
and the Jacobian condition number.


\paragraph{Definition.}

For species $s$, the adequacy certificate is
\begin{equation}\label{eq:err-adequacy}
\Delta_{\rm obs}^{(s),\text{bound}}
:= \frac{C(\xi_s, L_s)}
        {\sigma_{\min}(J_s)}\,
  D_{s,\geq 2}\,,
\end{equation}
where $C(\xi_s, L_s)$ is a constant that depends on the
quantum statistics and truncation level,
$\sigma_{\min}(J_s)$ is the smallest singular value of the
species Jacobian $J_s$
(Eq.~\ref{eq:chain-rule}), and $D_{s,\geq 2}$ is the
tangency diagnostic
(Eq.~\ref{eq:D-diagnostic}).  The certificate guarantees
\begin{equation}\label{eq:adequacy-bound}
\Delta_{\rm obs}^{(s),\text{actual}}
\leq \Delta_{\rm obs}^{(s),\text{bound}}
+ O\big(D_{s,\geq 2}^2\big)\,.
\end{equation}


\paragraph{Interpretation.}

The certificate is conservative: the bound overestimates
the actual error because it uses the worst-case direction
of the tangency residual.  The tightness ratio
\begin{equation}\label{eq:tightness}
\rho_s
:= \frac{\Delta_{\rm obs}^{(s),\text{actual}}}
        {\Delta_{\rm obs}^{(s),\text{bound}}}
\end{equation}
quantifies the conservatism.  For representative
configurations [CITATION NEEDED: Teff Paper~IV],
$\rho_s \sim 0.003$---the bound is never violated but is
three orders of magnitude loose.  The looseness is
acceptable for qualitative regime separation (identifying
when the $\Teff$ closure needs upgrading) but
insufficient for precision error propagation, which
requires the actual observable error (Level~5).


\paragraph{Scope limitation.}

The adequacy certificate is a \emph{local} bound: it
certifies the reconstruction quality at a single
conformal time $\eta$, not over the full line-of-sight
integration.  Cumulative error over cosmological
timescales is not bounded by the certificate
(Section~\ref{sec:limit-recovery}, ``Not established''
item~4).  The certificate should therefore be evaluated
at multiple epochs and combined with the observable error
(Level~5) for end-to-end accuracy assessment.


\paragraph{Phase~1.0 status.}

For neutrinos, $D_{\nu,\geq 2} = 0$ exactly
(free-streaming species have no collision source), so the
adequacy certificate is identically zero.  For photons
before recombination, $D_{\gamma,\geq 2}$ is nonzero but
suppressed by Thomson damping ($D_{\gamma,\geq 2} \sim
\sigma_{ab}/\dot{\tau} \sim 10^{-8}$), giving a certificate
$\Delta_{\rm obs}^{(\gamma)} \lesssim 10^{-8}$.  The
$\Teff$ closure is therefore certified adequate at all
epochs probed by the Phase~1.0 solver.


% ============================================================
\section{Level~4: Ladder-specific adequacy}
\label{sec:err-ladder}
% ============================================================

The fourth level tracks the error introduced by the
adaptive $\ell$-band routing of
Section~\ref{sec:ell-band}, which assigns different
computational strategies to different angular scales.


\paragraph{Definition.}

Three diagnostics measure the ladder-specific adequacy.
The deceleration-parameter adequacy is
\begin{equation}\label{eq:err-Dq}
\Delta_q^{(s)}
:= \left|
  q^{(s),\text{Teff}}_0
  - q^{(s),\text{ref}}_0
\right|\,,
\end{equation}
where $q_0$ is the present-epoch deceleration parameter
contributed by species $s$.  The number-to-energy ratio
adequacy is
\begin{equation}\label{eq:err-DNE}
\Delta_{N/E}^{(s)}
:= \left|
  \frac{j_s^a}{q_s^a}\bigg|_{\rm Teff}
  - \frac{j_s^a}{q_s^a}\bigg|_{\rm ref}
\right|\,,
\end{equation}
measuring the accuracy of the two-field reconstruction of
the number-to-energy flux ratio
(Proposition~\ref{prop:flux-independence}).  The moment
residual is
\begin{equation}\label{eq:err-Rmom}
\mathcal{R}_{A_k}^{(s),\text{mom}}
:= \left|
  m_s^{(k),\text{Teff}}
  - m_s^{(k),\text{ref}}
\right|\,,
\end{equation}
measuring the accuracy of the $k$-th observable moment
(Eq.~\ref{eq:obs-vector}).


\paragraph{Role in the adaptive strategy.}

These diagnostics are not post-hoc quality measures but
runtime decision variables.  In the Phase~4.0 adaptive
ladder (Section~\ref{sec:solver-programme}), the
$\ell$-band router evaluates $\Delta_q$ and
$\Delta_{N/E}$ at each epoch boundary and upgrades the
species representation (from one-field to two-field, or
from the $\Teff$ closure to the full PSTF hierarchy) when
the diagnostics exceed their tolerances.  The upgrade is
conservative by design: over-upgrading (using a more
expensive representation when a cheaper one would suffice)
wastes computational resources but does not introduce
error, while under-upgrading produces inaccurate results.
The frozen-rung time-slab strategy---upgrading only at
epoch boundaries rather than at every time step---avoids
the numerical instability that can arise from frequent
representation switching.


\paragraph{Phase~1.0 status.}

The Phase~1.0 solver does not implement adaptive $\ell$-band
routing; it operates at a fixed truncation
$\ell_{\rm max} = 2$ with a fixed species configuration
(one-field photons, two-moment neutrinos).  The
ladder-specific diagnostics are therefore not evaluated at
runtime.  The P1-3 discrepancy between the polarised and
unpolarised code paths ($3\%$ in $D_2$;
Section~\ref{sec:solver-performance}) is a surrogate
indicator: it measures the error introduced by the choice
of whether to include E-mode polarisation in the hierarchy,
which is the simplest form of a representation-switching
decision.


% ============================================================
\section{Level~5: Observable error}
\label{sec:err-observable}
% ============================================================

The fifth and final level measures the deviation of the
solver's output in the space of observables---the CMB power
spectra and spherical-harmonic coefficients---from the
reference solution.


\paragraph{Definition.}

The $a_{\ell m}$ error is
\begin{equation}\label{eq:err-alm}
\varepsilon_{a_{\ell m}}
:= \frac{
  \big|
    a_{\ell m}^{\rm Teff}
    - a_{\ell m}^{\rm ref}
  \big|
}{
  \big| a_{\ell m}^{\rm ref} \big|
}\,.
\end{equation}
The angular power spectrum error is
\begin{equation}\label{eq:err-Cl}
\varepsilon_{C_\ell}
:= \frac{
  \big|
    C_\ell^{\rm Teff}
    - C_\ell^{\rm ref}
  \big|
}{
  \big| C_\ell^{\rm ref} \big|
}\,.
\end{equation}
On anisotropic backgrounds, the off-diagonal covariance
$\langle a_{\ell m}\,a_{\ell' m'}^* \rangle$ is generically
nonzero (the Bianchi shear couples different $\ell$ and $m$),
and the covariance error
\begin{equation}\label{eq:err-cov}
\varepsilon_{\rm cov}
:= \frac{
  \big|
    \langle a_{\ell m}\,a_{\ell' m'}^* \rangle_{\rm Teff}
    - \langle a_{\ell m}\,a_{\ell' m'}^* \rangle_{\rm ref}
  \big|
}{
  \big|
    \langle a_{\ell m}\,a_{\ell' m'}^* \rangle_{\rm ref}
  \big|
}
\end{equation}
must be tracked for the direction-dependent programme
(Section~\ref{sec:phase30}).  The covariance error
captures the $\ell$--$\ell'$ mode mixing induced by
the background anisotropy, which is invisible to the
isotropic power spectrum $C_\ell$.


\paragraph{Relation to the Bayes factor.}

The Bayes factor $\ln \mathcal{B}$ computed by the
departure pipeline (Chapter~\ref{ch:results}) is a
functional of the likelihood, which in turn depends on
$D_2 \propto C_2$.  The observable error
$\varepsilon_{C_2}$ therefore propagates directly into the
evidence:
\begin{equation}\label{eq:err-lnB}
\delta(\ln \mathcal{B})
\sim \frac{\partial \ln \mathcal{B}}
          {\partial \ln D_2}\,
     \varepsilon_{C_2}\,.
\end{equation}
For the VER06 pipeline, the sensitivity
$\partial \ln \mathcal{B} / \partial \ln D_2 \sim
O(10)$, so an observable error of $\varepsilon_{C_2}
\sim 0.01$ ($1\%$) translates to
$\delta(\ln \mathcal{B}) \sim 0.1$---comparable to the
statistical uncertainty $\pm 0.07$ of the production
evidence.  The Phase~2.0 target
$\varepsilon_{C_2} < 0.02$ therefore ensures that the
systematic error from the solver is subdominant to the
statistical error of the inference.


\paragraph{Phase~1.0 status.}

The Phase~1.0 solver computes $D_2$ (proportional to
$C_2$ at $\ell = 2$) but does not compute the full
$C_\ell$ spectrum or $a_{\ell m}$.  The observable
error at $\ell = 2$ is $\varepsilon_{C_2} \simeq 0.08$
(the oracle deviation), which is the dominant error in
the current budget.  Phase~3.0 extends the observable
error to the full $a_{\ell m}$ spectrum, targeting
$\varepsilon_{a_{\ell m}} < 0.01$ at $\ell \leq 30$.


% ============================================================
\section{Cost--accuracy trade-off}
\label{sec:err-cost}
% ============================================================

The $\Teff$ programme's computational advantage lies in
selective upgrading: applying expensive representations
only where the diagnostics demand it, rather than
uniformly across all species and angular scales.  The
cost--accuracy trade-off quantifies this advantage.


\paragraph{Cost metrics.}

Two cost metrics are tracked alongside the error budget:
the wall-clock time per likelihood evaluation,
$\mathcal{C}_{\rm wall}$, and the memory footprint per
evaluation, $\mathcal{C}_{\rm mem}$.  For the Phase~1.0
solver, $\mathcal{C}_{\rm wall} = 420$~ms and
$\mathcal{C}_{\rm mem} \sim 50$~MB.  For comparison,
AniCLASS achieves $\mathcal{C}_{\rm wall} \sim 1$~s per
evaluation at $\ell_{\rm max} = 30$ in the linearised
regime.


\paragraph{Pareto frontier.}

The Pareto frontier in the
$(\mathcal{C}_{\rm wall}, \varepsilon_{\rm obs})$ plane
delineates the set of configurations that cannot be
improved in one dimension without degrading the other.
The $\Teff$ programme's adaptive strategy is designed to
operate near this frontier: the low-$\ell$ band uses
exact quadrature (high accuracy, moderate cost), the
mid-$\ell$ band uses the $\Teff$ ladder with selective
upgrading (controlled accuracy, reduced cost), and the
high-$\ell$ band uses the standard PSTF hierarchy
(baseline accuracy, baseline cost).  The selective-upgrading
result---that targeting upgrades where
$D_{s,\geq 2} > \varepsilon_{\rm thr}$ outperforms
uniform high-resolution computation in cost--accuracy
trade-off---is the operational vindication of the $\Teff$
approach.

For the scalar-amplitude pipeline, the Pareto frontier is
trivially achieved: only $\ell = 2$ matters, the
$\Teff$ closure is exact for the dominant neutrino
channel, and the evaluation cost ($420$~ms) is set by
the hierarchy ODE integration rather than by the
closure strategy.  The frontier becomes non-trivial for
the direction-dependent programme, where the full
$a_{\ell m}$ spectrum at $\ell \leq 30$ requires
$\sim\!10^3$ multipoles and the cost savings from
selective upgrading can reach an order of magnitude.


% ============================================================
\section{Phase~1.0 error budget summary}
\label{sec:err-phase10-budget}
% ============================================================

We summarise the error budget for the Phase~1.0 solver
and identify the dominant error at each level.

\begin{center}
\setlength{\tabcolsep}{5pt}
\begin{tabular}{llll}
\hline\hline
Level & Diagnostic & Phase~1.0 value & Dominant source \\
\hline
1. State-space
  & $\varepsilon_{\rm state}^{(\nu)}(\ell{=}2)$
  & $0.08$
  & Analytic $\Psi(k,\eta)$ \\[3pt]
2. Source-translation
  & $\varepsilon_{\rm Th}^{\rm lin}$
  & $\sim 10^{-3}$
  & Linearised Thomson \\[3pt]
3. Adequacy certificate
  & $\Delta_{\rm obs}^{(\gamma)}$
  & $\lesssim 10^{-8}$
  & Thomson suppression \\[3pt]
  & $\Delta_{\rm obs}^{(\nu)}$
  & $= 0$ (exact)
  & Free-streaming \\[3pt]
4. Ladder-specific
  & not evaluated
  & ---
  & Fixed truncation \\[3pt]
5. Observable
  & $\varepsilon_{C_2}$
  & $0.08$
  & $= \varepsilon_{\rm state}$ \\
\hline\hline
\end{tabular}
\end{center}

The dominant error is at Level~1 (state-space), arising
from the analytic gravitational potential.  The $\Teff$
closure itself contributes negligibly: the adequacy
certificate is $\leq 10^{-8}$ for photons and exactly
zero for neutrinos.  The source-translation error is
$\sim 10^{-3}$, three orders of magnitude below the
state-space error.  The error budget confirms that the
$\Teff$ approximation is not the bottleneck---the
gravitational potential is---and that Phase~2.0
(self-consistent Poisson) is the correct next step.

The error hierarchy also reveals a structural asymmetry
between species.  The neutrino sector contributes
$98.4\%$ of $D_2$ but introduces zero $\Teff$ error
(Levels~2--3), because free-streaming neutrinos reside
exactly on the $\Teff$ manifold.  The photon sector
contributes $1.6\%$ of $D_2$ and carries the only
nonzero $\Teff$ error, but this error is suppressed by
Thomson damping to a level that is irrelevant at any
foreseeable precision.  The $\Teff$ closure is therefore
not merely adequate but structurally optimal for the
species composition of the standard cosmological model:
the dominant species is exactly closed, and the
subdominant species is over-damped.

Before closing this chapter, we note that the five-level
hierarchy is designed to be extensible.  As the solver
programme advances through Phases~2.0--5.0
(Section~\ref{sec:solver-programme-ch9}), each level
acquires additional diagnostics---self-consistent
Poisson residuals at Level~1, nonlinear Thomson bridge
accuracy at Level~2, AniCLASS $a_{\ell m}$ comparison
at Level~5---without changing the overall structure.
The hierarchy is a framework for error accounting, not a
fixed set of numbers.


% ============================================================
\section{\texorpdfstring{\texttt{MioCertificate}}{MioCertificate}
         semantic rule}
\label{sec:err-mio-semantic}
% ============================================================

Sections~\ref{sec:err-state}--\ref{sec:err-phase10-budget}
treat error measurement as a numerical deviation between two
code paths.  The last three sections of this chapter shift
register: they clarify \emph{which kinds of claims the solver
ecosystem is permitted to make}, and they bind the low-level
error bounds of the five-level hierarchy to the architectural
constraints that prevent adequacy-at-one-level from being
re-labelled as adequacy-at-another.  The first of these
clarifications, presented here, is the semantic rule governing
the \texttt{MioCertificate} data object---the central
interface through which the Model-Independent Observatory
(MIO) reports departure diagnostics to downstream consumers.

The core claim is simple and load-bearing: a
\texttt{MioCertificate} is a \emph{diagnostic report}, not a
truth certificate.  The distinction is not rhetorical; it is
encoded in the dataclass schema, the constructor guards, and
the test battery that locks both.


\paragraph{Object definition.}

The \texttt{MioCertificate} (frozen dataclass defined in
\texttt{bass\_py/workspace/contracts/mio\_certificate.py})
carries the following fields:

\begin{itemize}
\item \emph{identification}: \texttt{report\_type},
  \texttt{probe\_name}, \texttt{channel};
\item \emph{diagnostic quantities}:
  \texttt{departure\_variables} (e.g.\ the median and
  $1\sigma$ width of $\Sigma^2_{\rm MIO}$),
  \texttt{adequacy\_indicators} (e.g.\ the
  $\ell$-independence $p$-value flag),
  \texttt{consistency\_metrics}
  (e.g.\ $\chi^2/\text{d.o.f.}$);
\item \emph{caveats}: \texttt{domain\_caveats},
  \texttt{channel\_caveats},
  \texttt{reduction\_status} $\in$ \{\texttt{theory-direct},
  \texttt{theory-approximate}, \texttt{diagnostic-only}\};
\item \emph{provenance}: \texttt{generated\_by},
  \texttt{git\_commit}, \texttt{config\_hash},
  \texttt{input\_data\_hashes};
\item \emph{cross-check hints}:
  \texttt{htt\_cross\_check\_suggested} (optional; agreement
  testing only, never evidence combination).
\end{itemize}

The schema contains \emph{no field} named
\texttt{posterior}, \texttt{posterior\_samples},
\texttt{log\_likelihood}, or any synonym.  This absence is
locked by a schema-hash digest test
(\texttt{test\_miocertificate\_schema\_frozen}).


\paragraph{Three forbidden interpretations.}

The MIO--HTT hard-separation rule of
Section~\ref{sec:err-g19} is enforced at the type layer by
three prohibitions on how a \texttt{MioCertificate} may be
consumed:

\begin{enumerate}
\item \emph{Not a posterior sample.} The method
  \texttt{MioCertificate.as\_posterior\_bundle()} is declared
  but intentionally raises \texttt{NotImplementedError} with
  a pointer to HTT's \texttt{PosteriorExportBundle}.  No
  schema field carries posterior semantics.
\item \emph{Not a truth certificate.} No downstream module is
  permitted to interpret the successful construction of a
  \texttt{MioCertificate} as an attestation that any
  particular cosmological model is correct.  The certificate
  reports \emph{what is directly visible in the data under
  model-independent reduction}, not \emph{what is true}.
\item \emph{Not an additive score.} A
  \texttt{MioCertificate} may not be combined with an HTT
  log-evidence, a matched-complexity score, or any other
  quantitative adequacy indicator via summation, averaging,
  or any reduction that produces a single merged master
  score.  Cross-checks (Section~\ref{sec:err-g19}) are
  explicitly non-additive.
\end{enumerate}


\paragraph{Test-layer anchors.}

Each prohibition is covered by a regression test under
\texttt{bass\_py/workspace/contracts/tests/}.  The field set
is locked by the schema-hash digest test so that any PR
attempting to add a \texttt{posterior} field breaks the hash
and fails in CI.  The ingestion-site refusal is tested by
\texttt{test\_htt\_cannot\_ingest\_miocertificate\_as\_likelihood}.
The additive-score ban is covered by a static-text regex scan
(\texttt{test\_g19\_enforcement}) that walks the production
subpackages and fails the build on any unreviewed hit.


\paragraph{The \texttt{reduction\_status} field.}

The three-valued \texttt{reduction\_status} makes the
adequacy claim of a certificate explicit:

\begin{itemize}
\item \texttt{theory-direct}---the reduction is a direct
  measurement under a defined model-independent observable.
  Examples: the directional-coherence Fisher $p$-value
  (Section~\ref{sec:mio-coherence}), the
  $\ell$-independence $\chi^2$ test.
\item \texttt{theory-approximate}---the reduction makes a
  controlled approximation (e.g.\ small-$\Sigma^2$
  expansion, flat-sky limit); the approximation and its
  domain of validity appear in
  \texttt{domain\_caveats}.
\item \texttt{diagnostic-only}---the reduction is too
  model-near to support a model-independent claim on its
  own but is retained as a diagnostic, typically behind a
  \texttt{production\_allowed=False} flag elsewhere.  The
  ZoA-20\textdegree{} default axis is the archetype: an
  anchor for diagnostic stability, not a production axis.
\end{itemize}

A \texttt{reduction\_status} of \texttt{diagnostic-only} is
a first-class signal that the certificate is \emph{not} to
be promoted into an inferential claim.  No reduction tier is
a truth certificate.


\paragraph{Why the schema is frozen.}

Freezing the schema via hash digest is not a stylistic
choice; it is the load-bearing defence against an
architectural drift in which a later PR adds a
\texttt{posterior} or \texttt{best\_fit\_params} field and
silently converts MIO output into a model-dependent
inference object.  The schema hash is tested in CI; any
drift fails loudly, and the failure message identifies the
field-name change rather than surfacing later as a silent
mis-labelling of a truth certificate in downstream code.


\paragraph{Historical correction.}

In an earlier draft of the research plan (v2), MIO was
described as a centralised post-hoc certification module; the
v2 vocabulary treated the \texttt{MioCertificate} as if it
issued a truth certificate.  The v3 reformulation, documented
in appendix A40 (G19 architectural stance), excises the v2
vocabulary and anchors adequacy reporting in the per-module
ownership table of Section~\ref{sec:err-epistemic}.  The
\texttt{legacy/} subtree preserves the v2 snapshot for
provenance with a prominent banner warning that its contents
must not be imported into active code paths.


% ============================================================
\section{Hard separation (G19) in practice}
\label{sec:err-g19}
% ============================================================

The second architectural rule enforced across the solver
ecosystem is the \emph{hard separation} between the
Model-Independent Observatory (MIO), the Bianchi-typed
Hypothesis Toolkit (HTT), and the Tangency--Stability--Chart
suite (TSC).  Referred to in the codebase as G19, it states:
\emph{the three adequacy domains never merge into a single
scalar score; they may only cross-check.}

This section records why the rule is non-negotiable, how it
is enforced, and which pipeline failure modes a relaxation
would introduce.


\paragraph{Statement.}

Under G19, the three code paths produce three independent
artefacts:

\begin{itemize}
\item HTT produces a posterior draw, a log-evidence, and a
  matched-complexity report;
\item TSC produces a realisability report, an admissibility
  Boolean, and the three-bound / Michaelis-Menten mirror of
  the bass Route-B solver;
\item MIO produces a \texttt{MioCertificate} with departure
  variables, adequacy indicators, caveats, and
  \texttt{reduction\_status}
  (Section~\ref{sec:err-mio-semantic}).
\end{itemize}

The artefacts may be compared pairwise, but a comparison
qualifies as a \emph{cross-check} only if it satisfies the
five-property definition of appendix A34: two
numerically-comparable scalars from two independently
implemented code paths; no merged field on the report
dataclass; a frozen \texttt{is\_cross\_check=True} tag; a
loud-fail guard; and no downstream consumption as a
likelihood term.  Objects that fail any of these conditions
are not cross-checks, and any object exposing a merged
estimator is a direct G19 violation.


\paragraph{What is forbidden.}

Three concrete operations are prohibited:

\begin{itemize}
\item summing or averaging an HTT log-evidence and a MIO
  departure statistic into a single master score;
\item relabelling a \texttt{MioCertificate} as an HTT
  likelihood term, posterior draw, or evidence
  contribution;
\item interpreting the successful construction of a
  \texttt{MioCertificate} as a truth certificate on behalf
  of any particular cosmological model.
\end{itemize}

Each prohibition corresponds to an explicit code-level
defence documented below.


\paragraph{Four-layer enforcement.}

G19 is enforced at four independent layers of the codebase,
so that no single human-review oversight defeats it:

\begin{itemize}
\item \textbf{Type layer.}  Frozen dataclass schemas forbid
  the offending field names at construction time.
  \texttt{MioCertificate} has no \texttt{posterior} field;
  \texttt{FFCrossCheckReport} carries
  \texttt{F\_Bayes\_tsc} and \texttt{F\_Bayes\_htt\_mean} as
  two distinct fields with no merged
  \texttt{F\_Bayes\_combined} counterpart; its
  \texttt{\_\_post\_init\_\_} hook raises
  \texttt{ValueError} if \texttt{is\_cross\_check} is
  negated.
\item \textbf{Generator-site layer.}  The
  \texttt{build\_mio\_certificate(\dots)} factory scans its
  keyword arguments for the substring \texttt{"posterior"}
  and raises before any instance is built, catching the case
  where a schema-allowed field is populated with a
  semantically forbidden value.
\item \textbf{Ingestion-site layer.}  HTT likelihood entry
  points type-check their inputs and refuse to accept a
  \texttt{MioCertificate}; the refusal is exercised by
  \texttt{test\_htt\_cannot\_ingest\_miocertificate\_as\_likelihood}.
\item \textbf{Repo-wide scan layer.}  A regex-based static
  scan (\texttt{test\_g19\_enforcement}) searches
  \texttt{bass\_py/\{bass, htt/htt, mio, src/common, tsc,
  workspace\}} for patterns such as
  \texttt{mio\_score + htt\_lnB} or
  \texttt{np.mean([mio\_score, htt\_lnB])}.  Any hit must be
  whitelisted explicitly or the build fails.
\end{itemize}


\paragraph{Currently wired cross-check channels.}

Five G19-compliant cross-check channels are presently wired
in the codebase (appendix A34):

\begin{itemize}
\item TSC-06 filling fraction
  (\texttt{tsc.integration.htt\_bridge} vs.
  \texttt{htt.core.analysis\_extended.FillingFraction}),
  agreeing at ${\rm rtol} < 10^{-6}$ on the S3 scenario;
  the only channel with a frozen
  \texttt{is\_cross\_check=True} dataclass
  (\texttt{FFCrossCheckReport});
\item TSC-03 three-bound hierarchy
  (\texttt{tsc.admissibility.three\_bound\_hierarchy}
  vs.\ \texttt{htt.core.bounds}), bit-identical at
  ${\rm rtol}=10^{-10}$ on the MES coefficient table;
\item TSC-05 Michaelis--Menten SSOT mirror
  (\texttt{tsc.charts.michaelis\_menten\_export} vs.\
  \texttt{bass.spectrum.cl\_assembly}), bit-identity
  (${\rm rtol}=0$);
\item PR13AH spherical-mean vs.
  \texttt{common.sky\_geometry} (\texttt{ChannelSummary}
  diff), the W4 R1 regression anchor;
\item MIO HJ-02a directional coherence vs.\ HTT dipole
  likelihood---advisory only, with no matched numerical
  pair, expressed through the certificate's
  \texttt{domain\_caveats}.
\end{itemize}

No merge channel exists; none is planned; any addition
would require a plan revision and is architecturally
forbidden (appendix A40).


\paragraph{Failure modes a relaxation would introduce.}

Relaxing G19 to allow merging produces four structural
failure modes (appendix A40):

\begin{enumerate}
\item \emph{Collapse of independent witnesses.}  A bug that
  corrupts both code paths identically becomes undetectable.
  TSC-06 exists precisely to produce two numerically
  comparable scalars from two independently implemented code
  paths; a merged score destroys that signal.
\item \emph{Miscounting of shared uncertainty.}  A merged
  filling-fraction does not reduce variance by $\sqrt 2$,
  because both paths share the same $(\epsilon_1,
  \epsilon_2, \epsilon_3)$ posterior draw.  Merged-variance
  arithmetic is therefore pseudo-statistical.
\item \emph{Contamination of HTT likelihoods.}  Once MIO and
  HTT scores are fungible, nothing prevents a later PR from
  ingesting a \texttt{MioCertificate} as an HTT likelihood
  input; the model-independence guarantee collapses and the
  certificate silently migrates into a truth certificate.
\item \emph{Destruction of the TSC realisability gate.}
  The realisability report is Boolean; a merged score forces
  the Boolean into an averaged continuous field and corrupts
  the admissibility gate.
\end{enumerate}

Each failure mode is architectural rather than quantitative:
the numerics might still pass the low-level error budget of
Sections~\ref{sec:err-state}--\ref{sec:err-observable} while
the resulting object ceases to be interpretable.  G19 is the
rule that keeps the error budget interpretable at all.


% ============================================================
\section{Distributed ownership of epistemic control}
\label{sec:err-epistemic}
% ============================================================

The third architectural clarification is that
epistemic-hygiene responsibility is \emph{distributed}: it
is not concentrated in any single module.  Each module in
the solver ecosystem owns its own adequacy and exposes that
ownership through a specific artefact; the ensemble of
artefacts never collapses to a single master score
(Section~\ref{sec:err-g19}).

This section records the ownership table and its
operational consequences for the error budget of this
chapter.


\paragraph{Principle.}

The distributed-ownership principle can be stated in a
single sentence: the solver tells you whether the numerics
are trustworthy; the admissibility module tells you whether
a configuration is physically realisable; the inference
module tells you whether a model's posterior passes
matched-complexity / LOOCV / PPC; the observatory tells
you whether the data is consistent with the model's
adequacy claims \emph{without referencing any specific
model}.  Each produces its own artefact and each is
consumed independently downstream.


\paragraph{Ownership table.}

The canonical ownership table (appendix A39) lists the
primary owner of each adequacy domain in the codebase; the
summary form is reproduced below.

\begin{center}
\setlength{\tabcolsep}{4pt}
\small
\begin{tabular}{lll}
\hline\hline
Domain & Primary owner & Contract artefact \\
\hline
Forward-solver fidelity
  & \texttt{bass.runtime.*}
  & \texttt{RuntimeGate} \\
Claim-tier labels
  & \texttt{bass.runtime.claims}
  & tier enum \\
Physical realisability
  & \texttt{tsc.admissibility.*}
  & \texttt{RealizabilityReport} \\
Tangency / entropy
  & \texttt{tsc.diagnostics.tangency}
  & invariant dict \\
Posterior adequacy
  & \texttt{htt.infer.matched\_complexity}
  & \texttt{MatchedComplexityReport} \\
Posterior predictive
  & \texttt{htt.advanced\_diagnostics.PosteriorPredictive}
  & \texttt{PPCReport} \\
Null-competition FPR
  & \texttt{htt.infer.null\_competition}
  & \texttt{NullCompetitionReport} \\
Model-indep. diagnostics
  & \texttt{mio.diagnostics.adequacy\_certificates}
  & \texttt{MioCertificate} \\
Directional coherence
  & \texttt{mio.coherence.directional}
  & \texttt{MioCertificate(directional\_coherence)} \\
Masked-sky caveats
  & \texttt{mio.diagnostics.masked\_sky\_caveats}
  & certificate \texttt{domain\_caveats} \\
Interface semantics
  & \texttt{bass\_py/workspace/contracts/}
  & 4 frozen dataclasses \\
Production-axis gate
  & \texttt{common.contracts.PreferredAxis}
  & \texttt{production\_allowed} Boolean \\
\hline\hline
\end{tabular}
\end{center}

Each row is independent; outputs across rows are never
summed.  No row occupies the role of a truth certificate.
The \texttt{MioCertificate} row ships a diagnostic report;
the \texttt{MatchedComplexityReport} row ships a
posterior-adequacy report under a specific model; the
\texttt{RealizabilityReport} row ships a Boolean
admissibility claim.  These three artefacts cross-check;
they do not merge.


\paragraph{Three first-line defences.}

Appendix A39 identifies three semantic first-line defences
that the ownership table relies on:

\begin{itemize}
\item \emph{Filename prefix.} MIO artefacts begin with
  \texttt{mio\_}
  (e.g.\ \texttt{mio\_directional\_coherence\_v1.json},
  \texttt{mio\_pr13am\_te\_sign\_v1.json}); enforced by
  \texttt{test\_reg02\_artifact\_prefix}.
\item \emph{Module tag.} MIO-owned modules physically
  located in other packages carry an
  \texttt{\_\_mio\_owned\_\_ = True} attribute plus a
  \texttt{\_\_mio\_rationale\_\_} string; the sole current
  example is
  \texttt{htt.PR13AM\_te\_sign\_d1d3\_bridge}.
\item \emph{Schema token ban.} The substring
  \texttt{"posterior"} is forbidden in any field name on
  \texttt{MioCertificate}, \texttt{AtlasEntry}, and the
  upstream \texttt{HttForwardOutput}; enforced by
  \texttt{test\_g19\_enforcement}.
\end{itemize}

Each defence is tested; together they prevent a later PR
from quietly re-concentrating adequacy reporting into a
single module---the category error that the v2 plan made
and that v3 corrects.


\paragraph{Historical correction.}

The v2 research plan described the Model-Independent
Observatory as a central module that issued a truth
certificate on behalf of the pipeline.  The v3 reformulation
(appendix A40) identifies that description as a category
error: adequacy is distributed, not centralised, and the MIO
output is a departure-variable report, not a truth
certificate.  Downstream code that was written against the
v2 draft must be migrated; the \texttt{legacy/} subtree
preserves the v2 snapshot with an explicit banner warning
against active use.


\paragraph{How the five-level hierarchy inherits the rule.}

The five-level error hierarchy of
Sections~\ref{sec:err-state}--\ref{sec:err-observable} is
itself an instance of distributed ownership.  Level~1
(state-space error) belongs to the forward solver; Level~2
(source translation) belongs to the $\Teff$ closure;
Level~3 (adequacy certificate) belongs to the tangency
diagnostic; Level~4 (ladder-specific) belongs to the
adaptive $\ell$-band router; Level~5 (observable) belongs
to the inference code.  No level consumes another level's
error as a truth certificate; each is diagnosed and
controlled on its own terms, and the Phase~1.0 budget
(Section~\ref{sec:err-phase10-budget}) is a table of
independent entries, not an additive sum.  The
architectural rules of
Sections~\ref{sec:err-mio-semantic}--\ref{sec:err-epistemic}
are therefore not auxiliary meta-commentary but the
logical closure of the error hierarchy itself.
